\documentclass[a4paper,11pt]{article}
\usepackage[a4paper,margin=1in]{geometry}
\usepackage{hyperref}
\usepackage{enumitem}

\title{Akash Kothapalli -- Resume Usage Guide}
\date{}

\begin{document}

\maketitle

This document describes how to use and maintain the different resume versions.

\section*{Resume Versions}

\begin{itemize}
\item \textbf{Backend Resume}
Used when applying for backend engineering roles.

Typical roles:
\begin{itemize}
\item Backend Developer
\item API Developer
\item Node.js Developer
\item Platform Engineer
\item Startup backend roles
\end{itemize}

\item \textbf{Full Stack Resume}
Used when applying for full-stack development roles.

Typical roles:
\begin{itemize}
\item Full Stack Developer
\item Web Developer
\item React + Node roles
\item Product engineering roles
\end{itemize}

\item \textbf{Software Engineer Resume}
Used when applying to large technology companies.

Examples:
\begin{itemize}
\item Amazon
\item Google
\item Microsoft
\item Adobe
\item Salesforce
\end{itemize}

\end{itemize}

\section*{Customization Rules}

Only the following sections should change between resumes:

\begin{itemize}
\item Profile section
\item Technical Skills order
\item Project descriptions
\end{itemize}

All other sections remain the same.

\section*{Application Strategy}

\begin{itemize}
\item Backend resume for startups and backend-heavy roles
\item Full-stack resume for product companies and web roles
\item Software engineer resume for big tech companies
\end{itemize}

\section*{Maintenance}

Whenever a new project or experience is added:

\begin{itemize}
\item Update the base resume
\item Propagate the changes to all three versions
\end{itemize}

\end{document}

\section*{Automatic Resume PDF Generation}

This repository uses GitHub Actions to automatically build PDF resumes.

Workflow:

\begin{itemize}
\item Update the resume \texttt{.tex} files in the \texttt{resumes/} directory
\item Commit and push the changes to GitHub
\item GitHub Actions automatically compiles the resumes
\item Generated PDF files can be downloaded from the \texttt{Actions} tab
\end{itemize}

No local LaTeX installation is required for generating PDFs.